% !TeX program = xelatex
% !TeX encoding = UTF-8
% Standalone source. Compile twice with XeLaTeX; no external figures or .bib.
% Project equations checked against CoupledBeams, commit
% 653173cd2a8b5f408832e1164c30a6651cbf80be (Version 0.5.1.8).
% This is a new theoretical note, not an implemented or tested damped solver.
\documentclass[11pt,a4paper]{article}
\usepackage[a4paper,left=23mm,right=23mm,top=22mm,bottom=23mm,headheight=15pt]{geometry}
\usepackage{fontspec}
\IfFontExistsTF{Liberation Serif}{\setmainfont{Liberation Serif}}{\setmainfont{Times New Roman}}
\IfFontExistsTF{Liberation Sans}{\setsansfont{Liberation Sans}}{\setsansfont{Arial}}
\IfFontExistsTF{DejaVu Sans Mono}{\setmonofont[Scale=0.82]{DejaVu Sans Mono}}{\setmonofont[Scale=0.85]{Courier New}}
\usepackage{polyglossia}
\setdefaultlanguage{russian}
\setotherlanguage{english}
\usepackage{amsmath,amssymb}
\usepackage{unicode-math}
\setmathfont{latinmodern-math.otf}
\usepackage{booktabs,tabularx,array}
\usepackage{enumitem}
\usepackage{microtype}
\usepackage{xcolor}
\usepackage{fancyhdr}
\usepackage{needspace}
\usepackage{tocloft}
\setlength{\cftbeforesecskip}{2pt}
\usepackage{hyperref}
\hypersetup{unicode=true,colorlinks=true,linkcolor=black,citecolor=black,urlcolor=blue!45!black,
 pdftitle={Поворотное соединение Кельвина—Фойгта: плоские колебания сопряжённых стержней},
 pdfsubject={Теоретическая постановка EB/RLB и определение комплексного спектра}}
\urlstyle{same}
\pagestyle{fancy}
\fancyhf{}
\fancyhead[L]{\small Поворотное соединение Кельвина—Фойгта}
\fancyhead[R]{\small EB / RLB}
\fancyfoot[C]{\thepage}
\renewcommand{\headrulewidth}{0.3pt}
\setlength{\parindent}{1.2em}
\setlength{\parskip}{0.22em}
\setlength{\emergencystretch}{2em}
\setlist{topsep=0.3em,itemsep=0.15em,parsep=0pt}
\setcounter{tocdepth}{1}
\numberwithin{equation}{section}
\setcounter{MaxMatrixCols}{12}
\allowdisplaybreaks[1]
\newcommand{\dd}{\,\mathrm d}
\newcommand{\ii}{\mathrm i}
\newcommand{\A}{\mathcal A}
\newcommand{\D}{\mathcal D}
\newcommand{\Sbeam}{\mathcal S}
\newcommand{\kth}{k_\theta}
\newcommand{\cth}{c_\theta}
\newcommand{\qJ}{q_J}
\newcommand{\TJ}{\mathcal T_J}
\newcommand{\KJ}{\mathcal K_J}
\newcommand{\BB}{\mathsf B}
\newcommand{\JJ}{\mathsf J}
\newcommand{\HH}{\mathsf H}
\newcommand{\TT}{\mathsf T}
\newcommand{\PP}{\mathsf P}
\newcommand{\CC}{\mathsf C_{\mathrm{cl}}}
\newcommand{\refsub}{\mathrm{ref}}
\newcommand{\tref}{t_{\refsub}}
\newcommand{\lref}{l_{\refsub}}
\newcommand{\Dref}{D_{\refsub}}
\newcommand{\mref}{m_{\refsub}}
\newcommand{\norm}[1]{\left\lVert#1\right\rVert}
\newcommand{\abs}[1]{\left|#1\right|}
\DeclareMathOperator{\diag}{diag}
\DeclareMathOperator{\rank}{rank}
\DeclareMathOperator{\tr}{tr}
\newcommand{\repo}[2]{\href{https://github.com/gimmiloot/CoupledBeams/blob/653173cd2a8b5f408832e1164c30a6651cbf80be/#1}{#2}}
\newcommand{\noteblock}[1]{\par\medskip\noindent\begin{minipage}{\linewidth}\hrule\smallskip#1\smallskip\hrule\end{minipage}\par\medskip}

\title{\vspace{-1.1em}\textbf{Поворотное соединение Кельвина—Фойгта\\при плоских колебаниях двух сопряжённых стержней}\\[0.5em]
\large Теоретическая постановка EB/RLB\\и определение комплексных собственных значений}
\author{\normalsize Рабочая теоретическая заметка проекта CoupledBeams}
\date{\normalsize 12 сентября 2026 г. \quad Версия заметки 1.0}

\begin{document}
\maketitle
\thispagestyle{empty}
\begin{abstract}
К ранее введённой поворотной пружине добавляется линейный вязкий демпфер,
действующий на относительную угловую скорость сечений в месте сопряжения.
Элементы соединены параллельно и образуют узел Кельвина—Фойгта. Стержни
остаются упругими; рассматриваются классическая модель Эйлера—Бернулли
с продольным движением (EB) и принятая в проекте симметричная слоистая
модель Reddy/FSDT (RLB). Из виртуальной работы и энергетического баланса
выводятся моментные условия. При временном множителе $e^{pt}$ задача
сводится к комплексному характеристическому уравнению
$\det\BB(p)=0$. Раздельно определяются частота затухающих колебаний
и показатель затухания. Обсуждаются продолжение собственных значений
от упругой задачи, локальное уточнение комплексного корня, контроль
невязок и обработка близких мод. Получены энергетические тождества
для собственных решений и формула первого приближения при слабом
демпфировании.
\end{abstract}

\noteblock{\textbf{Статус и происхождение.}
Координаты, усилия, энергии плеч и упругая сборка сохранены по документам
проекта на коммите \texttt{653173cd} (Version 0.5.1.8)~\cite{projectjoint,projectsource}.
Закон вязкого узла и его спектральное расширение выводятся в этой заметке;
они не приписываются уже выполненным упругим расчётам. Численный поиск
комплексного спектра и верификация вязкого решателя здесь не выполнялись.
Заметка является отдельным исходником, а не изменением файлов репозитория.}

\begingroup
\small
\setlength{\parskip}{0pt}
\tableofcontents
\endgroup
\newpage

\section{Исходная модель и обозначения}
\label{sec:model}

Рассматриваются два прямых стержня с постоянными вдоль каждого плеча
характеристиками. Их оси лежат в одной плоскости, внешние концы защемлены,
а внутренние соединены точечным безмассовым узлом. В каждом плече координата
$x_i\in[0,L_i]$ направлена \emph{от внешней заделки к узлу}; штрих означает
производную по собственной координате $x_i$, точка --- по времени.
Длины и характеристики плеч могут различаться.

Сохраняется физическое состояние проекта
\begin{equation}
 \symbfit y_i=(u_i,w_i,\psi_i,N_i,Q_i,M_i)^{\mathsf T}.
 \label{eq:state}
\end{equation}
Здесь $u_i,w_i$ --- локальные продольное и поперечное перемещения,
$\psi_i$ --- принятый поворот сечения, $N_i,Q_i,M_i$ --- продольная сила,
поперечная сила и изгибающий момент. Для локальных базисов
\begin{equation}
 \symbfit t_2=-c_\beta\symbfit t_1+s_\beta\symbfit n_1,
 \qquad
 \symbfit n_2=-s_\beta\symbfit t_1-c_\beta\symbfit n_1,
 \qquad c_\beta=\cos\beta,\quad s_\beta=\sin\beta.
 \label{eq:basis}
\end{equation}
Вектор поворота сечения и момент направлены соответственно как
$-\psi_i\symbfit k$ и $-M_i\symbfit k$, где одна и та же
$\symbfit k$ задаёт нормаль к плоскости. Эти знаки, как и порядок
\eqref{eq:state}, берутся из действующей постановки~\cite{projectjoint}.

Поступательная совместность в узле имеет вид
\begin{equation}
 u_1+c_\beta u_2+s_\beta w_2=0,
 \qquad
 w_1-s_\beta u_2+c_\beta w_2=0.
 \label{eq:translation}
\end{equation}
Во всех узловых равенствах значения относятся к $x_i=L_i$.
Угол $\beta$ характеризует исходную ненапряжённую конфигурацию и не является
дополнительной динамической координатой. В частности, моменты
$\kth\beta$ и $\cth\dot\beta$ не вводятся.

В RLB каждый стержень описывается пятью редуцированными характеристиками
\[
 \A_i>0,\quad \D_i>0,\quad \Sbeam_i>0,\quad m_i>0,\quad J_i>0.
\]
Они определяются прежней послойной интеграцией и редукцией при $B_i=0$,
$I_{1,i}=0$~\cite{projectsource}. Вязкоупругость самих слоёв в настоящей
заметке \emph{не вводится}: эти характеристики вещественны и не зависят
от частоты. $\A_i,\D_i,\Sbeam_i$ --- одномерные жёсткости, а не полные
матрицы ламината. $\Sbeam_i$ уже включает выбранную поправку на поперечный
сдвиг; её коэффициент не связан с параметрами демпфера.

\begin{center}
\begin{tabularx}{\linewidth}{@{}lXl@{}}
\toprule
Величина & Смысл & Единица \(\text{СИ}\)\\
\midrule
$\A_i,\Sbeam_i$ & Продольная и сдвиговая жёсткости & Н\\
$\D_i$ & Изгибная жёсткость & Н$\cdot$м$^2$\\
$m_i$ & Масса единицы длины & кг/м\\
$J_i$ & Вращательная инерция единицы длины & кг$\cdot$м\\
$\kth$ & Жёсткость поворотной пружины & Н$\cdot$м/рад\\
$\cth$ & Коэффициент поворотного вязкого сопротивления & Н$\cdot$м$\cdot$с/рад\\
$p$ & Комплексный временной показатель & с$^{-1}$\\
\bottomrule
\end{tabularx}
\end{center}
При анализе размерностей радиан считается безразмерным. Коэффициент
$\cth$ относится ко всему сосредоточенному узлу и не является вязкостью
материала, измеряемой в Па$\cdot$с.

\section{Как вводится демпфирование Кельвина—Фойгта}
\label{sec:kv}

Положим
\begin{equation}
 \qJ(t)=\Delta\psi(t)=\psi_1(L_1,t)-\psi_2(L_2,t).
 \label{eq:qj}
\end{equation}
Параллельное соединение означает, что пружина и демпфер испытывают
одинаковый относительный поворот, а передаваемые ими моменты складываются.
Именно эта схема называется элементом Кельвина—Фойгта~\cite{comsolkv}.
Идеальный поворотный демпфер создаёт момент, пропорциональный
\emph{относительной} угловой скорости его концов~\cite{mathworksdamper}.
Поэтому принимается закон
\begin{equation}
 \boxed{\TJ=\kth\qJ+\cth\dot\qJ,\qquad \kth\ge0,\quad \cth\ge0.}
 \label{eq:kv}
\end{equation}
В принятой системе знаков моментные условия на плечах будут
$M_1=-\TJ$ и $M_2=\TJ$; их вывод дан в разделе~\ref{sec:balance}.
Сопротивление абсолютному повороту каждого плеча относительно основания
описывалось бы другой физической схемой.

Упругая энергия узла и функция диссипации Рэлея равны
\begin{equation}
 U_J=\frac12\kth\qJ^2,
 \qquad
 \mathcal R_J=\frac12\cth\dot\qJ^{\,2}.
 \label{eq:UR}
\end{equation}
Функция $\mathcal R_J$ не является запасённой энергией. Мощность
необратимых потерь равна $2\mathcal R_J=\cth\dot\qJ^2$.
Использование этой функции не означает пропорциональную модель
демпфирования Рэлея $\mathbf C=a\mathbf M+b\mathbf K$.

При заданных независимых поворотах сечений обобщённые вязкие силы
выражаются матрицей
\begin{equation}
 \begin{bmatrix}Q^{\mathrm d}_{\psi_1}\\Q^{\mathrm d}_{\psi_2}\end{bmatrix}
 =-\cth
 \begin{bmatrix}1&-1\\-1&1\end{bmatrix}
 \begin{bmatrix}\dot\psi_1\\\dot\psi_2\end{bmatrix}.
 \label{eq:dampingmatrix}
\end{equation}
Матрица при $\cth>0$ положительно полуопределена и имеет ранг один.
Общий поворот с одинаковыми скоростями не приводит демпфер в действие.

Для постоянного относительного поворота вязкий момент равен нулю,
а статическая жёсткость остаётся $\kth$. При $\cth=0$ возвращается
прежняя упругая пружина. При $\kth=0$, $\cth>0$ узел передаёт вязкий
момент во время движения; это уже не свободный внутренний шарнир.
Идеальный шарнир получается только при $\kth=\cth=0$.

\section{Уравнения во времени, условия узла и энергия}
\label{sec:balance}

\subsection{Энергии и внутренние уравнения RLB}

Для плеч сохраняются энергии упругой RLB-постановки~\cite{projectjoint}:
\begin{align}
 U_i&=\frac12\int_0^{L_i}
 \left[\A_i(u_i')^2+\D_i(\psi_i')^2+
 \Sbeam_i(w_i'+\psi_i)^2\right]\dd x_i,\label{eq:U}\\
 T_i&=\frac12\int_0^{L_i}
 \left[m_i(\dot u_i^2+\dot w_i^2)+J_i\dot\psi_i^2\right]\dd x_i.
 \label{eq:T}
\end{align}
Внешние заделки задаются условиями $u_i(0)=w_i(0)=\psi_i(0)=0$.
Энергетически сопряжённые усилия
\begin{equation}
 N_i=\A_i u_i',\qquad
 Q_i=\Sbeam_i(w_i'+\psi_i),\qquad
 M_i=\D_i\psi_i'.
 \label{eq:constitutive}
\end{equation}
Обозначим $T=\sum_iT_i$, $U=\sum_iU_i$.
Принцип Лагранжа—д'Аламбера с вязкой виртуальной работой записывается как
\begin{equation}
 \delta\int_{t_0}^{t_1}(T-U-U_J)\dd t
 -\int_{t_0}^{t_1}\cth\dot\qJ\,\delta\qJ\dd t=0.
 \label{eq:dalembert}
\end{equation}
Вариации в начальный и конечный моменты времени равны нулю.
После интегрирования по времени и по $x_i$ объёмные члены дают
\begin{equation}
 \boxed{
 N_i'=m_i\ddot u_i,\qquad
 Q_i'=m_i\ddot w_i,\qquad
 M_i'=Q_i+J_i\ddot\psi_i.}
 \label{eq:timebalance}
\end{equation}
Вязкий член находится только в узле: распределённого демпфирования
плеч в \eqref{eq:timebalance} нет.

\subsection{Шесть условий сопряжения}

Граничная форма при узловых вариациях имеет вид
\begin{equation}
 \sum_{i=1}^2
 \bigl(N_i\delta u_i+Q_i\delta w_i+M_i\delta\psi_i\bigr)_{x_i=L_i}
 +(\kth\qJ+\cth\dot\qJ)(\delta\psi_1-\delta\psi_2)=0.
 \label{eq:boundaryvirtual}
\end{equation}
Оба узловых конца являются верхними пределами интегралов. Дополнительный
знак для второго плеча не вводится. Поступательные вариации связаны
условиями \eqref{eq:translation}, а поворотные вариации независимы.
Отсюда следуют равновесие сил и два моментных равенства
\[
 M_1+\kth\qJ+\cth\dot\qJ=0,
 \qquad
 M_2-\kth\qJ-\cth\dot\qJ=0.
\]
В прежнем порядке строк проекта полная система условий записывается как
\begin{equation}
 \boxed{\begin{aligned}
 u_1+c_\beta u_2+s_\beta w_2&=0,\\
 w_1-s_\beta u_2+c_\beta w_2&=0,\\
 M_1+\kth(\psi_1-\psi_2)+\cth(\dot\psi_1-\dot\psi_2)&=0,\\
 N_1-c_\beta N_2-s_\beta Q_2&=0,\\
 Q_1+s_\beta N_2-c_\beta Q_2&=0,\\
 M_1+M_2&=0.
 \end{aligned}}
 \label{eq:jointtime}
\end{equation}
Изменяется только закон в третьей строке. При $\cth=0$ система
совпадает с выведенными ранее условиями упругого узла~\cite{projectjoint}.

\subsection{Энергетический баланс и проверка знака}

Из \eqref{eq:constitutive} и \eqref{eq:timebalance} непосредственно следует
\[
 \frac{\dd}{\dd t}(T+U)
 =\sum_i[N_i\dot u_i+Q_i\dot w_i+M_i\dot\psi_i]_0^{L_i}.
\]
В заделках мощность равна нулю. Поступательные члены в узле взаимно
сокращаются благодаря совместности и равновесию сил. Моментные члены дают
\[
 M_1\dot\psi_1+M_2\dot\psi_2
 =-\TJ\dot\qJ
 =-\kth\qJ\dot\qJ-\cth\dot\qJ^2.
\]
Поэтому полная механическая энергия удовлетворяет тождеству
\begin{equation}
 \boxed{\dot E=-\cth\dot\qJ^2\le0,
 \qquad E=T+U+\tfrac12\kth\qJ^2.}
 \label{eq:energy}
\end{equation}
Положительное $\cth$ не может быть источником энергии. Однако
$\dot E\le0$ \emph{не доказывает затухания каждого движения до нуля}:
у решения с $\dot\qJ=0$ данный механизм потерь отсутствует.

\subsection{Независимая запись для классической EB}

Для EB
\begin{equation}
 U_i^{\mathrm{EB}}=\tfrac12\int_0^{L_i}
 [\A_i(u_i')^2+\D_i(w_i'')^2]\dd x_i,\qquad
 T_i^{\mathrm{EB}}=\tfrac12\int_0^{L_i}m_i(\dot u_i^2+\dot w_i^2)\dd x_i.
 \label{eq:ebenergies}
\end{equation}
Полагаются $\psi_i=-w_i'$, $M_i=-\D_iw_i''$,
$Q_i=-(\D_iw_i'')'$. Двукратное интегрирование изгибного члена
даёт ту же граничную форму $N_i\delta u_i+Q_i\delta w_i+M_i\delta\psi_i$.
Поэтому \eqref{eq:jointtime} и \eqref{eq:energy} остаются справедливыми.
Внутри каждого плеча
\begin{equation}
 (\A_i u_i')'=m_i\ddot u_i,
 \qquad
 (\D_i w_i'')''+m_i\ddot w_i=0.
 \label{eq:ebpde}
\end{equation}
В RLB подстановка $\psi=-w'$ без обнуления сдвиговой податливости
недопустима. Тем более нельзя заменить относительный поворот демпфера
разностью наклонов осевых линий при конечной $\Sbeam_i$.

\section{Комплексный спектральный параметр и частоты}
\label{sec:spectral}

\subsection{Временное соглашение}

Ищем одно собственное решение в виде
\begin{equation}
 \symbfit y_i(x_i,t)=\operatorname{Re}
 \{\widehat{\symbfit y}_i(x_i)e^{pt}\},
 \qquad p=-\alpha+\ii\omega_d.
 \label{eq:ansatz}
\end{equation}
Для колебательного решения выбирается представитель с $\omega_d>0$.
$\alpha=-\operatorname{Re}p$ --- показатель затухания амплитуды,
$\omega_d=\operatorname{Im}p$ --- угловая частота затухающих колебаний,
$f_d=\omega_d/(2\pi)$ --- частота в герцах. При $\alpha>0$
огибающая убывает как $e^{-\alpha t}$.

При вещественном отрицательном $p$ решение не колеблется и описывает
апериодическое затухание. Такое значение не нужно искусственно превращать
в положительную частоту или удалять из спектра. Для колебательной пары
удобны параметры
\begin{equation}
 \zeta=\frac{\alpha}{\sqrt{\alpha^2+\omega_d^2}},
 \qquad
 \delta_{\log}=\frac{2\pi\alpha}{\omega_d}.
 \label{eq:dampingratios}
\end{equation}
$\delta_{\log}$ относится к огибающей одного модального решения.
В общем непропорционально демпфированном случае величина $|p|$
не обязана совпадать с частотой соответствующей упругой моды.
Общая структура комплексного спектра механической системы второго
порядка обсуждается в~\cite{tisseur2001}.

Если вместо \eqref{eq:ansatz} используется $e^{\ii\widetilde\omega t}$,
то $\widetilde\omega=-\ii p=\omega_d+\ii\alpha$.
Следовательно, при этом соглашении \emph{положительная} мнимая часть
$\widetilde\omega$ соответствует затуханию. В настоящей заметке
основной неизвестной остаётся $p$; соглашения не смешиваются.

\subsection{Спектральный закон узла}

Из \eqref{eq:kv} и \eqref{eq:ansatz}
\begin{equation}
 \widehat{\TJ}=\KJ(p)\,\widehat{\qJ},
 \qquad \boxed{\KJ(p)=\kth+\cth p.}
 \label{eq:dynamicstiffness}
\end{equation}
Это динамическая жёсткость по относительному повороту, а не отношение
момента к угловой скорости. При $p=-\alpha+\ii\omega_d$
\[
 \KJ(p)=\kth-\cth\alpha+\ii\cth\omega_d.
\]
Поэтому при поиске собственных значений нельзя сохранять только
$\kth+\ii\cth\omega_d$ и отбрасывать $-\cth\alpha$.

При \emph{вынужденном} гармоническом движении с заданной вещественной
частотой $\omega$ допустима подстановка $p=\ii\omega$ и
$\KJ=\kth+\ii\omega\cth$. Но поиск максимумов частотного отклика
при вещественной $\omega$ --- другая задача. Он не заменяет определения
свободных собственных показателей \eqref{eq:ansatz}.

\subsection{Связь с нормировкой проекта}

Сохраняем фиксированные опорные величины
\begin{equation}
 \mref=\rho_0 A_{g,\refsub},\quad
 \Dref=E_0 I_{g,\refsub},\quad
 \tref=\lref^2\sqrt{\frac{\mref}{\Dref}},\quad
 z=p\tref.
 \label{eq:scales}
\end{equation}
$\tref$ имеет размерность времени. При нулевом демпфировании
$z=\ii\Omega$, где $\Omega=\omega\tref$ --- прежняя нормировка.
Введём независимые безразмерные параметры узла
\begin{equation}
 \boxed{\kappa_\theta=\frac{\kth\lref}{\Dref},\qquad
 d_\theta=\frac{\cth\lref}{\Dref\tref},\qquad
 \frac{\KJ(p)\lref}{\Dref}=\kappa_\theta+d_\theta z.}
 \label{eq:dimensionlessjoint}
\end{equation}
Такая нормировка пригодна и при $\kth=0$; параметр $\cth/\kth$
не нужен. Опорные величины не следует менять при вариации длины одного
плеча или текущих свойств ламината.

Для визуализации колебательной пары можно сохранять
\begin{equation}
 a=-\operatorname{Re}z=\alpha\tref,\quad
 \Omega_d=\operatorname{Im}z=\omega_d\tref,\quad
 \Lambda_d=\sqrt{\Omega_d}.
 \label{eq:lambdad}
\end{equation}
$\Lambda_d$ --- только вещественная характеристика частоты колебательной
части. Она не заменяет комплексного $z$ и ничего не говорит об $a$ без
отдельной записи. Для апериодических корней следует хранить сам $p$ и
показатель затухания, а не рисовать их как обычные частотные ветви.
Служебный параметр прежней проверки \texttt{epsilon\_limit} и отношение
толщины к длине не являются параметрами демпфирования.

\section{Точное характеристическое уравнение через перенос}
\label{sec:transfer}

\subsection{Матрица состояния с комплексным аргументом}

После подстановки \eqref{eq:ansatz} в \eqref{eq:timebalance} получаем
\begin{equation}
 \widehat{\symbfit y}_i'=\HH_i(p)\widehat{\symbfit y}_i,
 \qquad
 \HH_i^{\mathrm{RLB}}(p)=
 \begin{bmatrix}
 0&0&0&1/\A_i&0&0\\
 0&0&-1&0&1/\Sbeam_i&0\\
 0&0&0&0&0&1/\D_i\\
 m_i p^2&0&0&0&0&0\\
 0&m_i p^2&0&0&0&0\\
 0&0&J_i p^2&0&1&0
 \end{bmatrix}.
 \label{eq:Hp}
\end{equation}
Знаки инерционных членов здесь \emph{положительные перед $p^2$}.
При $p=\ii\omega$ они превращаются в прежние $-m_i\omega^2$,
$-J_i\omega^2$. Матрица классической EB равна
\begin{equation}
 \HH_i^{\mathrm{EB}}(p)
 =\HH_i^{\mathrm{RLB}}(p)\big|_{\,1/\Sbeam_i=0,\ J_i=0}.
 \label{eq:Heb}
\end{equation}
В пространственном состоянии по-прежнему шесть компонент на плечо.
Для элемента Кельвина—Фойгта не требуется отдельная внутренняя переменная
релаксации. Комплексные $p$ и амплитуды не являются добавлением физических
пространственных степеней свободы.

Поскольку коэффициенты постоянны по $x_i$, решение задаётся точно в рамках
одномерной модели:
\begin{equation}
 \TT_i(p,x_i)=\exp\bigl(\HH_i(p)x_i\bigr),\qquad
 \widehat{\symbfit y}_i(x_i)=\TT_i(p,x_i)\CC\symbfit r_i,
 \qquad
 \CC=\begin{bmatrix}0_{3\times3}\\I_3\end{bmatrix}.
 \label{eq:transfer}
\end{equation}
$\symbfit r_i=(N_i(0),Q_i(0),M_i(0))^{\mathsf T}$ --- комплексные
амплитуды реакций заделки. Для разных плеч обязательно используются
их собственные $L_i,\A_i,\D_i,\Sbeam_i,m_i,J_i$.

\subsection{Матрица узла и граничная матрица}

Опуская шляпки только внутри матричной записи, из \eqref{eq:jointtime}
получаем
\begin{equation}
 \JJ_{\mathrm{KV}}(p,\beta)=
 \begin{bmatrix}
 1&0&0&0&0&0&c_\beta&s_\beta&0&0&0&0\\
 0&1&0&0&0&0&-s_\beta&c_\beta&0&0&0&0\\
 0&0&\KJ&0&0&1&0&0&-\KJ&0&0&0\\
 0&0&0&1&0&0&0&0&0&-c_\beta&-s_\beta&0\\
 0&0&0&0&1&0&0&0&0&s_\beta&-c_\beta&0\\
 0&0&0&0&0&1&0&0&0&0&0&1
 \end{bmatrix},
 \label{eq:Jmatrix}
\end{equation}
где $\KJ=\kth+\cth p$. Матрица имеет размер $6\times12$ и ранг шесть
при любом комплексном $\KJ$: минор по столбцам
$(u_1,w_1,M_1,N_1,Q_1,M_2)$ равен единице. В частности, нуль
$\kth+\cth p$ не является особенностью исходной матрицы узла.

Введём
\begin{equation}
 \PP_i(p)=\TT_i(p,L_i)\CC,
 \qquad
 \boxed{\BB(p)=\JJ_{\mathrm{KV}}(p,\beta)
 \operatorname{blockdiag}\bigl(\PP_1(p),\PP_2(p)\bigr).}
 \label{eq:Bmatrix}
\end{equation}
Тогда для шести реакций
$\symbfit r=(\symbfit r_1^{\mathsf T},\symbfit r_2^{\mathsf T})^{\mathsf T}$
имеем
\begin{equation}
 \boxed{\BB(p)\symbfit r=0,\quad \symbfit r\ne0,
 \qquad f(p)=\det\BB(p)=0.}
 \label{eq:characteristic}
\end{equation}
Матрицы \eqref{eq:Hp} и \eqref{eq:Jmatrix} аналитичны по $p$;
матричная экспонента сохраняет аналитичность. Поэтому \eqref{eq:characteristic}
является нелинейной аналитической задачей на собственные значения
небольшой матрицы. Несмотря на размер $6\times6$, её спектр не ограничен
шестью корнями. Сам определитель в общем случае не является многочленом
второй или шестой степени по $p$. Этот тип нелинейных спектральных
задач рассматривается в~\cite{guettel2017}.

\subsection{Полезная структура одного поворотного канала}

Зависимость узловой матрицы от $\KJ$ находится только в третьей строке.
Линейность определителя по строке поэтому даёт точное представление
\begin{equation}
 f(p)=f_0(p)+(\kth+\cth p)f_1(p),\qquad f_r(-p)=f_r(p),\quad r=0,1.
 \label{eq:affine}
\end{equation}
Здесь $f_0$ --- определитель шарнирной сборки, а $f_1$ получается заменой
её третьей строки строкой относительного поворота, умноженной на карты
переноса; остальные строки сохраняются. Чётность $f_0,f_1$ следует из
зависимости плеч только от $p^2$.

Равенство \eqref{eq:affine} удобно для контроля реализации, но делить
его на $f_0$, $f_1$ или $\kth+\cth p$ без проверки нельзя. Так можно
потерять общие нули, в том числе движения с нулевым относительным
поворотом. Основным остаётся условие \eqref{eq:characteristic}.

При вещественных параметрах
\begin{equation}
 \BB(\overline p)=\overline{\BB(p)},\qquad
 f(\overline p)=\overline{f(p)}.
 \label{eq:conjugate}
\end{equation}
Корни образуют комплексно-сопряжённые пары или лежат на вещественной оси.
При $\cth>0$ симметрия $p\leftrightarrow-p$ в общем случае уже отсутствует.

\section{Вариационная спектральная задача и аналитические контроли}
\label{sec:weak}

Пусть $\varphi=(u_i,w_i,\psi_i)_{i=1,2}$ удовлетворяет заделкам и
поступательной совместности, а $v$ --- допустимое пробное поле.
Комплексные формы ниже линейны по первому аргументу и сопряжённо линейны
по второму:
\begin{align}
 \mathfrak m(\varphi,v)
 &=\sum_i\int_0^{L_i}
 [m_i(u_i\overline{v_{u,i}}+w_i\overline{v_{w,i}})
   +J_i\psi_i\overline{v_{\psi,i}}]\dd x_i,\label{eq:massform}\\
 \mathfrak a_0(\varphi,v)
 &=\sum_i\int_0^{L_i}
 [\A_i u_i'\overline{v_{u,i}'}+\D_i\psi_i'\overline{v_{\psi,i}'}
 +\Sbeam_i(w_i'+\psi_i)\overline{(v_{w,i}'+v_{\psi,i})}]\dd x_i,
 \label{eq:stiffform}\\
 \mathfrak a_k(\varphi,v)
 &=\mathfrak a_0(\varphi,v)+\kth q(\varphi)\overline{q(v)},
 \qquad q(\varphi)=\psi_1(L_1)-\psi_2(L_2).\label{eq:akform}
\end{align}
Для EB используются $\psi_i=-w_i'$, $J_i=0$,
изгибный член $\D_iw_i''\overline{v_{w,i}''}$ и отсутствует сдвиговой член.
В частности, в EB нет независимой массовой координаты $\psi$.

Из уравнений раздела~\ref{sec:balance} получаем
\begin{equation}
 \boxed{\mathfrak a_k(\varphi,v)
 +p\cth q(\varphi)\overline{q(v)}
 +p^2\mathfrak m(\varphi,v)=0\quad\text{для всех }v.}
 \label{eq:weakpencil}
\end{equation}
Это квадратичный \emph{операторный} пучок. После точного исключения
пространственной зависимости он приводит к неполиномиальной матрице
\eqref{eq:Bmatrix}; противоречия здесь нет.

\subsection{Отсутствие растущих собственных движений}

Положим $v=\varphi$ и введём вещественные величины
\[
 \mathcal M_\varphi=\mathfrak m(\varphi,\varphi)>0,\qquad
 \mathcal K_\varphi=\mathfrak a_k(\varphi,\varphi)>0,\qquad
 \mathcal C_\varphi=\cth\abs{q(\varphi)}^2\ge0.
\]
Положительность упругой формы обеспечивается положительными жёсткостями
и внешними заделками: нулевая деформация каждого плеча вместе с заделкой
даёт нулевое поле. Для ненулевой физической собственной формы масса
положительна. Из \eqref{eq:weakpencil}
\begin{equation}
 p^2\mathcal M_\varphi+p\mathcal C_\varphi+\mathcal K_\varphi=0.
 \label{eq:modalidentity}
\end{equation}
При $p=\sigma+\ii\nu$, $\nu\ne0$, мнимая часть равна
$\nu(2\sigma\mathcal M_\varphi+\mathcal C_\varphi)=0$. Следовательно,
\begin{equation}
 \boxed{\alpha=-\sigma=
 \frac{\cth\abs{q(\varphi)}^2}{2\mathcal M_\varphi},\qquad
 \omega_d^2=\frac{\mathcal K_\varphi}{\mathcal M_\varphi}-\alpha^2.}
 \label{eq:exactalpha}
\end{equation}
Эти равенства точны для \emph{уже найденной комплексной собственной формы},
а не только для слабого демпфирования. Но подставить в них произвольную
упругую форму и получить точный вязкий спектр нельзя: $\varphi$ сама
зависит от $p$ и $\cth$.

Для вещественного $p\ge0$ левая часть \eqref{eq:modalidentity} строго
положительна. Поэтому такой корень невозможен. В итоге все собственные
показатели рассматриваемой пассивной модели удовлетворяют
$\operatorname{Re}p\le0$, а $p=0$ не является собственным значением.
Это спектральное утверждение и энергетическое тождество; здесь не
доказываются полнота модального разложения и равномерное экспоненциальное
затухание всех решений распределённой системы.

\subsection{Движения, не чувствующие демпфера}

Если упругая собственная форма $\varphi_0$ при частоте $\omega_0>0$
удовлетворяет $q(\varphi_0)=0$, то при $p=\pm\ii\omega_0$ оба узловых
слагаемых с $\kth$ и $\cth$ в \eqref{eq:weakpencil} исчезают для любого $v$.
Значит, та же форма остаётся собственным незатухающим движением при
любых конечных неотрицательных $\kth,\cth$, если остальные параметры
не меняются. Обратно, из \eqref{eq:exactalpha} при $\cth>0$ следует:
у колебательного собственного движения с $\alpha=0$ относительный
поворот обязательно равен нулю.

Для одинаковых плеч отражение проекта имеет вид
\[
 \mathscr P(\symbfit y_1,\symbfit y_2)
 =(F\symbfit y_2,F\symbfit y_1),\qquad
 F=\diag(1,-1,-1,1,-1,-1).
\]
Оно сохраняет RLB- и EB-уравнения и квадратичные узловые формы.
В классе $\symbfit y_2=\eta F\symbfit y_1$, $\eta=-1$,
имеем $\psi_2=\psi_1$, поэтому весь класс не демпфируется этим элементом.
В классе $\eta=+1$ моментное условие половинной задачи равно
$M+2(\kth+\cth p)\psi=0$. При неодинаковых плечах такое разделение
в общем случае неприменимо: нужны полная матрица и непривязанные
к прежним классам формы. Упругие основания этого разграничения
зафиксированы в~\cite{projectrobustness}.

Малый, но ненулевой $q$ не даёт точной защиты от потерь. Кроме того,
один положительный коэффициент $\cth$ не обеспечивает одинакового
затухания всех мод и не доказывает монотонного роста их показателей
затухания при неограниченном увеличении $\cth$.

\section{Слабое демпфирование и взаимодействие близких мод}
\label{sec:small}

\subsection{Первая поправка к простому собственному значению}

Пусть $\omega_{0j}$ --- простая изолированная частота упругой задачи
с уже заданной $\kth$, а $\varphi_{0j}$ --- её форма.
Для вывода введём \emph{только здесь} малый безразмерный множитель
$\epsilon_d$: $\cth=\epsilon_d c_*$, $c_*>0$ фиксировано.
Он не связан ни с геометрической тонкостью, ни с прежним
\texttt{epsilon\_limit}.

Ищем $p=p_0+\epsilon_dp_1+O(\epsilon_d^2)$,
$\varphi=\varphi_0+\epsilon_d\varphi_1+O(\epsilon_d^2)$,
где $p_0=\pm\ii\omega_{0j}$. Дифференцирование
\eqref{eq:weakpencil} при $\epsilon_d=0$ и подстановка
$v=\varphi_{0j}$ сокращают члены с $\varphi_1$ благодаря
самосопряжённости упругой задачи. Получаем
\[
 2p_0p_1\mathcal M_{0j}+p_0c_*\abs{q(\varphi_{0j})}^2=0,
 \qquad \mathcal M_{0j}=\mathfrak m(\varphi_{0j},\varphi_{0j}).
\]
Отсюда
\begin{equation}
 \boxed{
 p_j^\pm(\epsilon_d)=\pm\ii\omega_{0j}
 -\epsilon_d\frac{c_*\abs{q(\varphi_{0j})}^2}{2\mathcal M_{0j}}
 +O(\epsilon_d^2).}
 \label{eq:smallp}
\end{equation}
В первом порядке изменяется вещественная часть $p$, а изменение
частоты $\omega_d$ начинается не ранее второго порядка для такого
простого изолированного корня. Формула годится как начальное приближение
и независимый аналитический контроль локального продолжения, но не как
замена решения \eqref{eq:characteristic} при произвольном $\cth$.

Связь с уже использовавшейся упругой чувствительностью
\begin{equation}
 s_j=\frac{(\Dref/\lref)\abs{q(\varphi_{0j})}^2}
 {\omega_{0j}^2\mathcal M_{0j}}
 =\frac{\partial\log\omega_{0j}^2}{\partial\kappa_\theta}
 \label{eq:sensitivity}
\end{equation}
даёт при малом $d_\theta$ и $\Omega_{0j}=\omega_{0j}\tref$
\begin{equation}
 \alpha_j\tref\simeq\tfrac12d_\theta\Omega_{0j}^{\,2}s_j,
 \qquad
 \zeta_j\simeq\tfrac12d_\theta\Omega_{0j}s_j.
 \label{eq:smallnondim}
\end{equation}
Чувствительность $s_j$ сама по себе не равна коэффициенту затухания.
Утверждение \eqref{eq:smallnondim} относится к слабому возмущению
изолированной моды; численные свойства вязкой системы из прежних
упругих графиков непосредственно не следуют.

\subsection{Почему нельзя демпфировать каждую упругую моду независимо}

Пусть вещественные упругие формы при фиксированной $\kth$ нормированы
по массе: $\mathfrak m(\varphi_{0r},\varphi_{0s})=\delta_{rs}$.
При конечном приближении
$\varphi\approx\sum_{j=1}^{n} a_j\varphi_{0j}$ положим
$b_j=q(\varphi_{0j})$. Проекция \eqref{eq:weakpencil} даёт
\begin{equation}
 \left[p^2 I_n+p\cth\symbfit b\symbfit b^{\mathsf T}
 +\diag(\omega_{01}^2,\ldots,\omega_{0n}^2)\right]\symbfit a=0.
 \label{eq:modalqep}
\end{equation}
Матрица демпфирования имеет ранг не более одного, но обычно содержит
ненулевые \emph{внедиагональные} элементы. Поэтому назначить каждой форме
свой скалярный коэффициент потерь и отбросить взаимодействие ---
дополнительное приближение, особенно опасное около veering.
Различие пропорционального и непропорционального демпфирования
обсуждается в~\cite{tisseur2001}.

Для близкой пары следует хотя бы сохранять весь блок
\begin{equation}
 \begin{bmatrix}
 p^2+\omega_{0a}^2+p\cth b_a^2 & p\cth b_ab_b\\
 p\cth b_ab_b & p^2+\omega_{0b}^2+p\cth b_b^2
 \end{bmatrix}.
 \label{eq:twomode}
\end{equation}
Его корни служат объяснением и возможными предикторами; конечное число
упругих форм не превращает этот блок в точное решение распределённой
задачи. Итоговые собственные значения проверяются по полной
\eqref{eq:characteristic}. В точке кратности нужно рассматривать
собственное подпространство, а не применять \eqref{eq:smallp} к случайному
SVD-базису. Нулевой относительный поворот определённой комбинации
кратных упругих форм при этом остаётся точным основанием отсутствия потерь.

\section{Как практически искать комплексные собственные значения}
\label{sec:numerics}

\subsection{Что сохраняется и что меняется в прежнем методе}

Сохраняются физическое состояние, экспоненты переноса, сборка по узлу
и восстановление формы из реакций. Меняется область поиска: уравнение
\eqref{eq:characteristic} решается по двум вещественным компонентам $p$
либо непосредственно в комплексной арифметике.

Поиск смен знака определителя на вещественной оси, применявшийся к
упругим частотам, теперь неприменим. У комплексного числа нет такого
порядка знаков. Минимум $\sigma_{\min}(\BB(\ii\omega))$ при вещественной
$\omega$ тоже не гарантирует собственного значения: затухающий полюс
обычно находится вне мнимой оси. Минимизация $|f|$ без проверки нулевой
невязки может остановиться в ненулевом минимуме.

В физическую матрицу входят $p^2$ и $\kth+\cth p$, а не $|p|^2$,
$\operatorname{Im}p$ или $\overline p$. Нельзя просто разрешить комплексную
жёсткость в старом вещественном интерфейсе, оставив вещественную проверку
аргумента, приведение амплитуд к \texttt{float} или старые инерционные члены.
Это вопросы будущей реализации, а не изменения уже выполненного кода.

\subsection{Масштабирование и локальная коррекция}

До решения вводятся согласованные масштабы перемещений, реакций и строк.
Например,
\[
 U_*=\lref,\qquad M_*=\Dref/\lref,\qquad F_*=\Dref/\lref^2.
\]
Для реакций и строк соответственно положим
\begin{align*}
S_a&=\diag(F_*,F_*,M_*,F_*,F_*,M_*),\\
S_r&=\diag(U_*^{-1},U_*^{-1},M_*^{-1},F_*^{-1},F_*^{-1},M_*^{-1}).
\end{align*}
Безразмерный вектор $\symbfit a=S_a^{-1}\symbfit r$
удовлетворяет задаче с безразмерной матрицей
\begin{equation}
 \widehat{\BB}(z;d_\theta)=S_r\BB(z/\tref;\cth)S_a
 \label{eq:scaledB}
\end{equation}
Обратимые масштабы не меняют собственных значений.
При вычислении аналитической производной масштабы фиксируются.
Нормы, модули, SVD и адаптивное выравнивание могут применяться в
диагностике, но не должны незаметно входить в функцию, для которой
используется комплексная аналитическая производная.

Для простого корня удобна локальная система на собственное значение
и вектор реакций:
\begin{equation}
 \widehat{\BB}(z)\symbfit a=0,\qquad
 \symbfit h^*\symbfit a=1,
 \label{eq:augmented}
\end{equation}
где $\symbfit h$ фиксирован во время коррекции и выбран так, чтобы
нормирующее произведение не обращалось в ноль. Шаг Ньютона имеет вид
\begin{equation}
 \begin{bmatrix}
 \widehat{\BB}(z)&\widehat{\BB}_z(z)\symbfit a\\
 \symbfit h^*&0
 \end{bmatrix}
 \begin{bmatrix}\Delta\symbfit a\\\Delta z\end{bmatrix}
 =-
 \begin{bmatrix}\widehat{\BB}(z)\symbfit a\\
 \symbfit h^*\symbfit a-1\end{bmatrix}.
 \label{eq:newton}
\end{equation}
Это система размера $7\times7$ в комплексных неизвестных, эквивалентная
системе из 14 вещественных уравнений. Она не требует минимизировать
модуль определителя. Общие методы Ньютона для нелинейных собственных
задач рассматриваются в~\cite{guettel2017}; \eqref{eq:newton} получено
линеаризацией конкретной системы \eqref{eq:augmented}.

В качестве альтернативы допускается решение
$\operatorname{Re}f(p)=\operatorname{Im}f(p)=0$ с согласованным
вещественным Якобианом. В обоих случаях требуется близкое начальное
приближение и последующая проверка полной матрицы. Нуль определителя
сам по себе не восстанавливает идентичность моды.

\subsection{Производная передаточной матрицы}

При необходимости производная вычисляется без предположения о
коммутативности:
\begin{align}
 \partial_p\TT_i(p,L_i)
 &=\int_0^{L_i}e^{\HH_i(p)(L_i-s)}\,
       \partial_p\HH_i(p)\,e^{\HH_i(p)s}\dd s,\label{eq:derivativeT}\\
 \partial_p\BB
 &=\partial_p\JJ_{\mathrm{KV}}\,\operatorname{blockdiag}(\PP_1,\PP_2)
 +\JJ_{\mathrm{KV}}\operatorname{blockdiag}(\partial_p\PP_1,\partial_p\PP_2).
 \label{eq:derivativeB}
\end{align}
В $\partial_p\HH_i$ ненулевые элементы равны $2m_ip$ на местах $(4,1)$,
$(5,2)$ и $2J_ip$ на месте $(6,3)$. В $\partial_p\JJ_{\mathrm{KV}}$
ненулевы только $\cth$ и $-\cth$ в поворотных столбцах третьей строки.
Формулу \eqref{eq:derivativeT} можно вычислять как верхний правый блок
экспоненты
\[
 \exp\begin{bmatrix}
 L_i\HH_i&L_i\partial_p\HH_i\\0&L_i\HH_i
 \end{bmatrix}.
\]
Упрощение $\partial_p e^{L\HH}=Le^{L\HH}\HH_p$ без доказанного
$[\HH,\HH_p]=0$ неверно. В безразмерной коррекции
$\partial_z\BB=\partial_p\BB/\tref$ при фиксированном $\cth$.

\subsection{Ограниченное продолжение от упругого спектра}

Для первого вычислительного этапа естественна следующая последовательность.
\begin{enumerate}
\item Зафиксировать геометрию, характеристики плеч и $\kappa_\theta$.
      Взять проверенные упругие корни и формы той же постановки;
      начальные показатели $z_j(0)=\ii\Omega_{0j}$.
\item При малом положительном $d_\theta$ использовать
      \eqref{eq:smallp} или \eqref{eq:smallnondim} только как предиктор.
      Для близкого или кратного набора начинать с подпространства;
      при точной симметрии допустимо заранее разделить независимые классы.
\item Уточнять каждую новую пару $(z,\symbfit a)$ по
      \eqref{eq:newton} или эквивалентной комплексной характеристической
      задаче. Сохранять фактический корень и восстановленную форму.
\item Продолжать по $d_\theta$ с соседними решениями как предикторами.
      Шаг уменьшать адресно при плохой коррекции, близких полюсах
      или неоднозначной форме, а не глобально во всём диапазоне.
\item Сопоставлять \emph{комплексные} формы через массовую корреляцию
      с комплексным сопряжением. Для RLB учитывать $J_i|\psi_i|^2$.
      Различать происхождение от упругой моды и сортировку по $\omega_d$.
\item Комплексно-сопряжённую пару учитывать как одну колебательную
      моду реального движения. Если пара подходит к вещественной оси,
      сохранять обе дальнейшие вещественные ветви при их появлении;
      это два показателя апериодического отклика, а не пропущенные частоты.
\end{enumerate}

Комплексные собственные формы в общем случае не ортогональны по обычной
массовой метрике. Не нужно принудительно ортогонализовать разные
разрешённые собственные решения и выдавать полученные комбинации за моды.
При совпадении собственных значений возможны как независимые формы,
так и дефектный корень с присоединёнными векторами~\cite{tisseur2001}.
Поэтому малый MAC или вырождение корректора не следует автоматически
«исправлять» перестановкой ближайших частот.

Такое продолжение отслеживает выбранных потомков упругого спектра.
Оно не является доказательством, что найдены все комплексные корни
в некоторой области. Прежнее понятие «седьмого guard» на вещественной
частотной оси нельзя механически перенести на комплексную плоскость.

\subsection{Контроль полноты только при конкретной необходимости}

Если нужно выяснить число корней внутри отдельной подозрительной области,
теоретически применим принцип аргумента:
\begin{equation}
 N_\Gamma=\frac{1}{2\pi\ii}\oint_\Gamma
 \tr\bigl(\BB(p)^{-1}\BB_p(p)\bigr)\dd p.
 \label{eq:argument}
\end{equation}
Для аналитической исходной матрицы без полюсов и контура без собственных
значений это число нулей с алгебраической кратностью. Численная реализация
требует контроля квадратуры и невырожденности матрицы на контуре;
контурные методы описаны в~\cite{guettel2017}. Равенство приведено как
адресный способ проверки кластера, а не как требование выполнять дорогой
контурный аудит для каждой обычной точки продолжения.

\subsection{Что проверять у каждого принятого корня}

В безразмерных согласованных переменных вычисляются
\begin{equation}
 r_B=\frac{\norm{\widehat{\BB}\symbfit a}}
 {\norm{\widehat{\BB}}\norm{\symbfit a}},
 \qquad
 r_\sigma=\frac{\sigma_{\min}(\widehat{\BB})}
 {\sigma_{\max}(\widehat{\BB})}.
 \label{eq:residualB}
\end{equation}
Затем в физических переменных проверяются заделки, все шесть условий
\eqref{eq:jointtime} с $\partial_t\mapsto p$ и восстановление обоих плеч.
Момент должен получаться из решения плеча, а не присваиваться по формуле
пружины перед проверкой этой же формулы.

Дополнительный энергетический остаток имеет вид
\begin{equation}
 r_E=\frac{|p^2\mathcal M_\varphi+p\mathcal C_\varphi+\mathcal K_\varphi|}
 {|p|^2\mathcal M_\varphi+|p|\mathcal C_\varphi+\mathcal K_\varphi}.
 \label{eq:residualE}
\end{equation}
Он полезен вместе с \eqref{eq:exactalpha}, но не заменяет граничных невязок.
Проверяются сопряжённость \eqref{eq:conjugate}, отсутствие объяснимо
неподдержанного роста $\operatorname{Re}p>0$ и возврат к упругим решениям.

Малость $r_B$ не является безусловной оценкой ошибки $p$, особенно у
кратного или плохо обусловленного корня. Для простого корня чувствительность
коррекции зависит от величины
$\symbfit w^*\widehat{\BB}_z\symbfit a$, где $\symbfit w$
--- левый нуль-вектор. Если она мала, нельзя обещать много верных знаков
на основании одного остатка. Допуски для $\omega_d$, малого $\alpha$ и
физических условий задаются раздельно \emph{до} расчёта. Неразрешимый
показатель затухания нельзя выдавать за измеренное число, а теоретически
нулевое затухание нельзя заменять искусственным отрицательным
$\operatorname{Re}p$ ради «устойчивой» картинки.

\section{Предельные случаи и минимальные проверки реализации}
\label{sec:limits}

\paragraph{Нулевая вязкость.}
При $\cth=0$ должны совпасть прежние упругая матрица при вещественной
$\omega$ и новая матрица при $p=\ii\omega$, включая физические знаки.
Показатели равны $p_j=\pm\ii\omega_{0j}$. Достаточна проверка в уже
сохранённых совместимых корнях; полный повтор упругих карт не требуется.

\paragraph{Переход RLB в EB.}
При $1/\Sbeam_i\to0$, $J_i\to0$ и фиксированных
$\A_i,\D_i,m_i,L_i,\beta,\kth,\cth$ матрица \eqref{eq:Hp} переходит в
\eqref{eq:Heb}. При каждом фиксированном конечном комплексном $p$
непрерывность экспоненты даёт тот же переход для $\TT_i,\PP_i,\BB$.
Одного выключения сдвиговой податливости при сохранённом $J_i$
недостаточно для классической EB. Из матричного предела не следует
равномерная сходимость всего бесконечного спектра.

\paragraph{Нулевая жёсткость и статический случай.}
При $\kth=\cth=0$ возвращается моментно-свободный шарнир с передачей сил.
При $p=0$ вязкость не влияет на статические условия. Для данной конструкции
с внешними заделками отсутствие нулевой собственной формы следует из
положительности энергии, а не из исключения нуля поисковым алгоритмом.

\paragraph{Жёсткий узел и большая вязкость.}
При $\kth\to\infty$, фиксированных $p$ и $\cth$, деление моментной строки
на $\kth$ даёт $\Delta\psi\to0$ для ограниченных состояний; точный
жёсткий узел по-прежнему следует задавать отдельно. Большая $\cth$
не является универсальной заменой жёсткого узла во всём спектре:
предел при фиксированном $p\ne0$ не описывает автоматически ветви,
у которых $p$ само стремится к нулю или уходит из рассматриваемой области.

\paragraph{Неактивный демпфер.}
У доказанно сохраняющейся формы с $q=0$ должны остаться
$p=\pm\ii\omega_0$ и нулевая мощность потерь. Это корректный контроль,
а не отказ от требования затухания. В одинаковой конструкции соответствующий
класс сохраняется в обеих теориях; в неодинаковой нельзя создавать его
проекцией решения на старую симметрию.

\paragraph{Скалярный контроль временного соглашения.}
Для отдельной тестовой системы с одной поворотной инерцией $I_*>0$
\[
 I_*\ddot q+\cth\dot q+\kth q=0,
 \qquad
 p_\pm=\frac{-\cth\pm\sqrt{\cth^2-4I_*\kth}}{2I_*}.
\]
Она проверяет знаки, сопряжённость и переход от колебательного к
апериодическому движению. Это тест алгоритма на известном уравнении,
не замена безмассового узла исходной конструкции инерционным элементом
и не независимая верификация стержневой модели.

\section{Итог постановки и границы заметки}

Новая физическая гипотеза состоит в одном локальном законе
\eqref{eq:kv}. Она не изменяет линейно-упругие свойства плеч и добавляет
неотрицательную диссипацию только в относительном повороте.
Для поиска собственных колебаний требуется находить комплексное $p$
из \eqref{eq:characteristic}, используя \eqref{eq:Hp} и
$\KJ(p)=\kth+\cth p$. После этого отдельно определяются
$\omega_d$, $\alpha$ и, при необходимости, $\Lambda_d$.

Для первого численного этапа предлагается ограниченное продолжение
нескольких уже проверенных упругих мод по $d_\theta$ при фиксированных
геометрии и $\kappa_\theta$, с полной комплексной коррекцией и физическим
контролем каждого нового решения. Модальные приближения
\eqref{eq:smallp}--\eqref{eq:twomode} служат предикторами и средствами
объяснения, но не подменяют итоговые корни передаточной задачи.

Настоящий вывод не определяет численное значение $\cth$ для реального
изделия и не доказывает применимость идеального линейного демпфера на
произвольных частотах, скоростях и температурах. Не учитываются сухое
трение, люфт, инерция и эксцентриситет узла, вязкоупругость слоёв,
несимметричный ламинат с $B\ne0$ и внешнее возбуждение. Свойства
затухания конкретных EB/RLB-конфигураций предстоит проверить отдельным
вычислительным этапом. Сохранённые ограничения исходной RLB-верификации
новым теоретическим выводом не снимаются.

\appendix
\section{Конечномерная квадратичная задача как вспомогательное представление}
\label{app:qep}

После выбора конечного числа допустимых координат или упругих форм
вариационная система имеет вид
\[
 (p^2\mathbf M+p\mathbf C+\mathbf K)\symbfit a=0.
\]
Её можно линеаризовать, введя $\symbfit v=(\symbfit a,p\symbfit a)^{\mathsf T}$:
\begin{equation}
 \begin{bmatrix}0&I\\-\mathbf K&-\mathbf C\end{bmatrix}\symbfit v
 =p\begin{bmatrix}I&0\\0&\mathbf M\end{bmatrix}\symbfit v.
 \label{eq:linearization}
\end{equation}
Это стандартная связь квадратичной и обобщённой линейной собственных
задач~\cite{tisseur2001}. Для \eqref{eq:modalqep} используются
$\mathbf M=I$, $\mathbf C=\cth\symbfit b\symbfit b^{\mathsf T}$,
$\mathbf K=\diag\omega_{0j}^2$.

При конечном модальном усечении получаются приближённые показатели;
близость к точному спектру требует проверки по размеру подпространства
или по исходной передаточной матрице. Матрица \eqref{eq:Bmatrix} после
точного переноса не является матрицей вида
$p^2\mathbf M+p\mathbf C+\mathbf K$ с постоянными шестимерными коэффициентами.
Поэтому \eqref{eq:linearization} не даёт автоматического способа
линеаризовать её без дополнительной аппроксимации. В данной работе
такое усечение и новый решатель не реализуются.

\section{Основания и воспроизведение документа}

Упругие энергии, знаки усилий, локальные координаты, порядок состояния
и структура сборки взяты из документов проекта~\cite{projectjoint,projectsource}.
Описание параллельной схемы и идеального вязкого поворотного элемента
опирается на~\cite{comsolkv,mathworksdamper}. Переход к \eqref{eq:jointtime},
комплексная матрица \eqref{eq:Hp}, энергетические тождества
\eqref{eq:energy}, \eqref{eq:exactalpha} и прикладная схема продолжения
выведены здесь для данной конструкции. Общие спектральные методы
сопоставлены с~\cite{tisseur2001,guettel2017}, а не приписаны старым
расчётным отчётам.

Список источников встроен в файл; внешние рисунки, база \texttt{.bib}
и данные расчётов для компиляции не требуются. Исходник рассчитан на
XeLaTeX. Обычно достаточно двух запусков:
\begin{quote}\small\ttfamily
xelatex inplane\_kelvin\_voigt\_joint\_theory.tex\\
xelatex inplane\_kelvin\_voigt\_joint\_theory.tex
\end{quote}
Если компилятор просит повторить сборку для ссылок, нужен ещё один запуск.
При подготовке проверяются компиляция и отображение формул; это не
численная верификация вязкой механической модели. Предполагаемое место
исходника при последующем добавлении в проект ---
\path{docs/laminated_beams/inplane_kelvin_voigt_joint_theory.tex}.

\begingroup
\raggedright
\renewcommand{\refname}{Источники}
\begin{thebibliography}{9}
\addcontentsline{toc}{section}{Источники}
\bibitem{projectjoint}
Проект CoupledBeams.
\emph{Поворотно-упругий узел двух стержней при плоских колебаниях}.
Теоретическая заметка и последующие уточнения.
\repo{docs/laminated_beams/inplane_rotational_spring_joint.md}{\nolinkurl{docs/laminated_beams/inplane_rotational_spring_joint.md}}.
Версия основания: commit \texttt{653173cd2a8b5f408832e1164c30a6651cbf80be},
Version 0.5.1.8.

\bibitem{projectsource}
Проект CoupledBeams.
\emph{Source contract: Reddy Chapter 4, symmetric laminated beam}.
\repo{docs/laminated_beams/reddy_ch4_source_contract.md}{\nolinkurl{docs/laminated_beams/reddy_ch4_source_contract.md}}.
Та же версия репозитория. Используется документированная проектная
редукция; реконструкция источника не приравнивается к полной
независимой верификации RLB.

\bibitem{projectrobustness}
Проект CoupledBeams.
\emph{Поворотная пружина: слоистая RLB и небольшое различие длин}.
\repo{docs/laminated_beams/inplane_spring_robustness.md}{\nolinkurl{docs/laminated_beams/inplane_spring_robustness.md}}.
Та же версия репозитория. Основание для упругих контрольных конфигураций,
не источник результатов вязкого расчёта.

\bibitem{comsolkv}
COMSOL Multiphysics.
\emph{Linear Viscoelasticity}, раздел \emph{The Kelvin--Voigt Model}.
Structural Mechanics Module User's Guide, версия 6.4.
\href{https://doc.comsol.com/6.4/doc/com.comsol.help.sme/sme_ug_theory.06.029.html}{Официальная документация COMSOL}.
Обращение: 12.09.2026.

\bibitem{mathworksdamper}
MathWorks. \emph{Rotational Damper: Viscous damper in mechanical rotational systems}.
Simscape Documentation.
\href{https://www.mathworks.com/help/simscape/ref/rotationaldamper.html}{Официальное описание идеального поворотного демпфера}.
Обращение: 12.09.2026.

\bibitem{tisseur2001}
Tisseur F., Meerbergen K.
The quadratic eigenvalue problem.
\emph{SIAM Review}. 2001. Vol.~43, No.~2. P.~235--286.
\href{https://doi.org/10.1137/S0036144500381988}{DOI: 10.1137/S0036144500381988}.
\href{https://eprints.maths.manchester.ac.uk/466/1/38198.pdf}{Открытая авторская копия}.

\bibitem{guettel2017}
Güttel S., Tisseur F.
The nonlinear eigenvalue problem.
\emph{Acta Numerica}. 2017. Vol.~26. P.~1--94.
\href{https://doi.org/10.1017/S0962492917000034}{DOI: 10.1017/S0962492917000034}.
\href{https://eprints.maths.manchester.ac.uk/2538/1/main.pdf}{Открытая авторская копия}.
\end{thebibliography}
\endgroup
\end{document}
